\chapter{Über den Satz von Bolzano}

\epigraf{Cuanto peor, mejor para todos, y cuanto peor para todos, mejor, mejor para mí el suyo, beneficio político.}{Mariano Rajoy}

\lettrine[loversize=0.2, lines=2]{É}{s} una llei no escrita de la realitat que, entre dues condicions contràries, n’hi ha d’haver una tercera que no és cap de les dues i que, tanmateix, participa de les dues. Aquesta llei no és una observació empírica, sinó una necessitat formal. No prové de l’experiència sinó de la condició mateixa del canvi.

Imagina un entorn on dues formes de força —dues estratègies vitals— es troben condicionades per un únic paràmetre. No una pluralitat de factors, no la complexitat del món real, sinó una sola variable que determina, sense excepció, la prevalença de l’una o de l’altra. És així com es pot concebre una situació teòrica, però no per això menys eloqüent, en què un lleó i un tauró lluiten.

Aquest combat no és una metàfora, sinó un sistema. El lleó, adaptat a la sequedat i al moviment sobre sòl ferm, domina quan l’entorn és àrid. El tauró, en canvi, gràcies a la seva massa i estabilitat, guanya terreny quan l’aigua puja. El paràmetre és clar: el nivell del líquid. A mesura que augmenta, la força relacional es desplaça. En un extrem, un és superior. En l’altre, ho és l’altre.

A mesura que pugem lentament el nivell de l’aigua, el lleó, tot i mantenir el seu instint depredador, perd gradualment eficàcia, mentre que el tauró la guanya. No hi ha un punt sobtat en què tot canviï; hi ha, en canvi, un procés lent i subtil en què les forces de les dues bèsties s’acosten lentament fins a trobar-se en una mena d’indefinició, un instant en què cap de les dues no pot afirmar amb claredat el seu domini. Podríem acceptar, doncs, que existeix un cert nivell d'aigua on les dues forces de la natura s'igualen i cap dels dos contendents guanya el combat.

Tanmateix, aquesta observació porta a una pregunta més inquietant. Acceptant la gradualitat del procés, és legítim preguntar-nos fins a quin punt aquesta continuïtat és realment absoluta. Imaginem, per un moment, la disminució gradual del nivell d'aigua fins a arribar al punt exacte en què l'equilibri es trenca. Quin és el grau mínim d'aigua que separa la victòria inevitable del tauró de la del lleó? Què succeeix quan extraiem tan sols una gota? I què succeeix, més radicalment encara, quan reduïm aquesta gota a l'escala de les molècules?

En aquest punt, l'observador es veu forçat a admetre la paradoxa inherent al concepte mateix de canvi continu. Podem imaginar clarament les dues situacions extremes i oposades, però ens resulta gairebé impossible concebre amb claredat el moment exacte del canvi. Aquesta dificultat no és accidental ni un simple error perceptiu; és més aviat un reflex de la naturalesa fonamental del món.

I aquí emergeix, inevitablement, el dubte sobre la validesa d'aquest continu absolut. Potser la realitat, vista a prou distància, ens sembla suau i contínua, però, a mesura que ens hi apropem, descobrim que hi ha salts discrets, fronteres microscòpiques, línies invisibles que la nostra ment no pot travessar sense un esforç conceptual extrem. Si una sola molècula d'aigua pot decidir la victòria del tauró o del lleó, aleshores, en última instància, el món revela la seva naturalesa no pas com una transició gradual, sinó com una successió subtil de ruptures infinitesimals.

I potser, en última instància, hem d'acceptar que aquest punt intermedi, aquest instant d'equilibri absolut, no existeix realment sinó com una mena d'abstracció necessària, una frontera conceptual creada per la ment humana. Una frontera que, vista des d'una escala prou petita, es dissol en el mar d'incerteses microscòpiques on cada molècula i cada instant quàntic semblen qüestionar incessantment la nostra visió contínua i ordenada del món.

\par\bigskip
Ara traslladem aquesta sospita —que entre dos extrems inevitables habita
forçosament un punt intermedi— a un escenari encara més abismal:  
imagina que deixem caure una bomba termonuclear al mig d’un desert
perfectament pla.  

A distàncies molt curtes de l’hipocentre\footnote{%
El punt zero d’una explosió nuclear rep, en llenguatge tècnic, el nom
d’\emph{hipocentre}. També s'utilitza per a terratrèmols.}, qualsevol matèria s’esvaeix
en plasma; a distàncies molt grans, la sorra ni tan sols arriba a
enrojolar‑se.  Hi ha, tanmateix, un anell invisible —una franja d’unes
poques desenes de metres— on la temperatura màxima instantània
oscil·la entorn dels \(350\ \si{\celsius}\): prou calor per fondre el
mozzarella, daurar la crosta i cuinar el llevat, però insuficient per
transformar‑ho tot en carbonissa.

En aquest annulus, la pizza assoleix la cocció perfecta; ni crua ni
cremada, sinó exactament al límit de l’exquisit.  Entre el regne de
l’infracuït i el regne del calcinat hi ha, inevitablement, la latitud
del punt just.  I si la temperatura decau de manera contínua amb la
distància —premissa que l’atmosfera respecta fora de turbulències
locals—, aleshores el \emph{Satz von Bolzano} ens xiuxiueja que aquesta
franja ha d’existir: no és un caprici gastronòmic, sinó una conseqüència
formal.

Però torna la pregunta: què passa quan un sol gra de sorra altera
l’albedo local, quan un petit corrent ascendent afegeix turbulència,
quan la discontinuïtat quàntica d’un fotó desvia l’energia d’una vora
de la crosta?  La línia perfecta del gust es fragmenta en una pols
d’atzars microscòpics.  Bolzano ens regala la certesa abstracta
d’aquesta corona ideal; la natura, amb la seva exquisida rugositat,
ens recorda que el continu no és mai indemne al minúscul, i que
dins de cada teorema batega, com una broma còsmica, la impugnació
de les seves pròpies hipòtesis.

Així, fins i tot la pizza termonuclear revela la mateixa lògica que
governava el duel entre el lleó i el tauró: la certesa d’un punt
intermedi, i a la vegada la sospita íntima que, si escrutem prou
endins, aquest punt es dissoldrà en miríades de fronteres
impossibles de fixar.  Continuïtat proclamada; discontinuïtat
subterrània.  Heus aquí, una vegada més, la ironia fonamental que
fa ballar la realitat al compàs del nostre afany de mesura.

\par\bigskip
Així, arribem a una curiosa dualitat que sembla travessar tota situació en què governa la continuïtat aparent: d'una banda, tenim un tauró i un lleó en lluita, dos éssers de carn i ossos separats per una frontera d'aigua; de l'altra, una pizza i una bomba termonuclear, objectes inerts separats per una frontera de temperatura. Dos escenaris aparentment allunyats, però units per la mateixa paradoxa matemàtica. Dues cares d'una única moneda: el món dels sistemes vius i el món de la matèria inanimada revelant-nos, junts, l'eterna dansa del discret i el continu.